\section{Synthetic experiment details}
\label{app:synthetic}

This appendix records the settings, artifacts and complete tables behind the
synthetic measurements quoted in the main text. Every number below is taken
from a named file under \texttt{results/}; where a value is a mean, a median
or a fitted slope that we computed from the raw rows of such a file rather
than being printed in it, we say so at that place.

\subsection{Generator, architectures and common protocol}
\label{app:synthetic_setup}

The generator produces hyperspectral images $X \in \R^{H \times W \times S}$
with $H = W = 16$ and $S = 64$ and a two-class per-pixel label. A known
spectral direction $u \in \R^S$ is placed in each pixel, scaled by a per-pixel
latent $z \sim \mathcal{N}(0,1)$ and modulated by a scalar strength $\alpha$;
smooth per-class spatial templates, modulated by a scalar strength $\beta$,
enter the label-generating logits. At $\alpha = \beta = 1$ and noise level
$0.1$ the Bayes-optimal accuracy is approximately $0.85$--$0.95$.

The spectral module $f_\thetap$ is a linear per-pixel map $\R^{64} \to \R^{16}$
without bias ($\Cf = 1024$ throughout). Four spatial modules are used across
the experiments: \emph{SpatialMLP}, a per-pixel one-hidden-layer network of
width $M$ (no spatial receptive field, $\Cg = \Theta(M)$); \emph{SpatialDeepMLP},
three hidden layers all $M \times M$ ($\Cg = \Theta(M^2)$); \emph{SpatialCNN},
two $3 \times 3$ convolutional layers; and \emph{SpatialViT}, a two-layer
pre-LN pixel-token transformer. The per-pixel loss is softmax cross-entropy
averaged over sites.

Two initializations of the spatial module are distinguished throughout, and
the spectral module's initialization is never altered. \emph{Default} is
PyTorch's \texttt{nn.Linear} initialization, whose bias variance is
$\Theta(1/\mathrm{fan\text{-}in})$. \emph{Theorem} initialization draws weights
$\mathcal{N}(0, 2/\mathrm{fan\text{-}in})$ and biases $\mathcal{N}(0,
\sigma_b^2)$ with $\sigma_b = 0.5$ held fixed in width, which is what the
curvature-disparity theorem's head hypothesis requires. Three further
controls decouple the weight law from the bias law:
\texttt{default\_fixedbias} (default weights, fixed-variance biases),
\texttt{he\_scaledbias} (He weights, Gaussian biases at the default variance
$1/(3\,\mathrm{fan\text{-}in})$) and \texttt{he\_zerobias} (He weights, zero
biases).

\subsection{Experiment 1.1: curvature blocks at initialization}
\label{app:exp11}

\paragraph{What is measured.} For each (architecture, width, seed) we compute
$\lambda_{\max}$ of the two diagonal Gauss--Newton blocks $G_{\thetap\thetap}$
and $G_{\phip\phip}$ of the averaged cross-entropy loss at random
initialization, in float64, on one fixed batch of 64 images per seed, by
Lanczos with full reorthogonalization applied to a matrix-free GGN block
operator ($80$ Lanczos steps, top-1). Before any sweep the runner executes a
validation gate that cross-checks the operator's $\lambda_{\max}$ against a
densely materialized $J^\top H_\tau J$ eigendecomposition on tiny
configurations for both architectures and both blocks, and aborts unless the
relative error is below $10^{-6}$. This closes the loss-Hessian versus
Gauss--Newton measurement gap present in the first version of this experiment:
the quantity measured here is the Gauss--Newton block the theory is stated for.
Five seeds per width. Artifacts:
\texttt{results/exp1\_1\_v3\_ggn\_aggregated.csv},
\texttt{results/exp1\_1\_v3\_ggn\_theoreminit\_aggregated.csv} and
\texttt{results/exp1\_1\_v3\_ggn\_sigmab\_isolation\_aggregated.csv}, each with
per-width seed means and standard deviations.

\paragraph{Results.} Table~\ref{tab:app_exp11} lists the two shallow-head arms
and the five deep-head arms. Eigenvalue entries are the per-width means over
five seeds printed in those files; the log--log slopes are ordinary
least-squares fits of $\log \lambda$ on $\log M$ over the full width range,
computed by us from those per-width means.

\begin{table}[htbp]
\centering
\small
\setlength{\tabcolsep}{4pt}
\caption{Experiment 1.1 (v3): top Gauss--Newton block eigenvalues at
initialization, five seeds per width. Each eigenvalue is given at the smallest
and the largest width of the arm's range. Slopes are least-squares log--log
fits over that range, computed by us from the per-width seed means in the three
aggregated CSVs; they use three orders of magnitude of width for the shallow
head and two for the deep head, and are finite-range descriptions rather than
asymptotic exponents. The five deep-head rows are the five arms of the
initialization isolation.}
\label{tab:app_exp11}
\begin{tabular}{lrrrrrrr}
\toprule
& & \multicolumn{3}{c}{$\lambda_\thetap$} & \multicolumn{3}{c}{$\lambda_\phip$} \\
\cmidrule(lr){3-5}\cmidrule(lr){6-8}
Initialization & $M$ range & first & last & slope & first & last & slope \\
\midrule
\multicolumn{8}{l}{\emph{Shallow head} (one hidden layer, $\Cg = \Theta(M)$)} \\
default & $16$--$8192$ & $0.0184$ & $0.0220$ & $-0.019$ & $0.62$ & $46.0$ & $0.705$ \\
theorem & $16$--$8192$ & $0.691$ & $0.538$ & $-0.040$ & $2.14$ & $441.6$ & $0.895$ \\
\midrule
\multicolumn{8}{l}{\emph{Deep head} (three hidden $M\times M$ layers, $\Cg = \Theta(M^2)$)} \\
default & $16$--$1024$ & $0.00032$ & $0.00037$ & $+0.012$ & $0.68$ & $1.11$ & $0.105$ \\
theorem & $16$--$1024$ & $0.371$ & $0.398$ & $-0.006$ & $3.62$ & $243.1$ & $0.989$ \\
default, fixed bias & $16$--$1024$ & $0.00055$ & $0.00063$ & $-0.041$ & $1.67$ & $73.8$ & $0.924$ \\
He, scaled bias & $16$--$1024$ & $0.310$ & $0.343$ & $+0.030$ & $1.50$ & $27.6$ & $0.696$ \\
He, zero bias & $16$--$1024$ & $0.336$ & $0.372$ & $-0.050$ & $1.45$ & $17.6$ & $0.606$ \\
\bottomrule
\end{tabular}
\end{table}

\paragraph{Reading.} The spectral cap is flat in width in every arm: the seven
$\lambda_\thetap$ slopes lie in $[-0.050, +0.030]$, and the level shifts by
three orders of magnitude between initializations without acquiring a width
trend. On the shallow head under the theorem's initialization the spatial
block grows with slope $0.895$ from $2.14$ to $441.6$; under default fan-in
initialization the same architecture gives slope $0.705$. The deep head
separates the two mechanisms most sharply: slope $0.105$ under default
initialization against $0.989$ under theorem initialization. The isolation
arms show that the width-linear behaviour survives if \emph{either} the
weight-variance term or the bias-variance term of the head hypothesis is
width-independent --- fixed-variance biases with default weights reach slope
$0.924$, He weights with default-variance biases $0.696$ and He weights with
zero biases $0.606$ --- while default fan-in initialization, which makes both
terms vanish with width, does not produce it. The corresponding curvature
ratios $\Dcurv$, first and last width, are $39.8 \to 2405$ and $3.4 \to 869$
for the two shallow arms and $3167 \to 3060$, $18.1 \to 668$, $5552 \to
1.17\times10^{5}$, $5.7 \to 84.5$ and $4.6 \to 46.7$ for the five deep arms, in
table order. They illustrate why the slope and the level must be quoted
together, and why we do not use $\Dcurv$ to attribute differences across arms:
the deep default arm has the largest ratio of any arm and the smallest
$\lambda_\phip$ slope, because its $\lambda_\thetap$ is $10^{-4}$ rather than
$10^{-1}$.

\subsection{Experiment 1.2 (v4): joint and frozen trajectories}
\label{app:exp12v4}

Data as in Appendix~\ref{app:synthetic_setup} with $\alpha = \beta = 1.0$ and
noise $0.10$; $512$ training and $128$ test images; Adam at learning rate
$10^{-3}$, batch size $32$, $150$ epochs ($2400$ steps); three seeds
($42$,$43$,$44$). The frozen arm sets \texttt{requires\_grad} to false on the
spectral parameters, so the spectral gradient is identically zero there by
construction. Widths are architecture-specific: $D \in \{16, 64, 256, 1024\}$
for the CNN head and $D \in \{64, 128, 256\}$ for the ViT head.
Table~\ref{tab:app_exp12v4} gives the full grid; means and standard deviations
over the three seeds were computed by us from the per-run rows of
\texttt{results/exp1\_2v4\_summary.csv}.

\begin{table}[htbp]
\centering
\small
\setlength{\tabcolsep}{4pt}
\caption{Experiment 1.2 (v4): endpoint and peak test accuracy and endpoint
losses, mean $\pm$ sd over three seeds, computed from the per-run rows of
\texttt{results/exp1\_2v4\_summary.csv}. \emph{Peak} is the maximum over the
run's per-epoch test evaluations and is therefore selection-optimistic.
$\Cg$ is the spatial module's parameter count; $\Cf = 1024$ in every row.}
\label{tab:app_exp12v4}
\begin{tabular}{llrrcccc}
\toprule
Head & Arm & $D$ & $\Cg$ & Final test acc. & Peak test acc. & Final test loss & Final train loss \\
\midrule
CNN & joint  & $16$   & $2{,}610$     & $0.618 \pm 0.061$ & $0.682 \pm 0.055$ & $0.650 \pm 0.037$ & $0.453 \pm 0.002$ \\
CNN & frozen & $16$   & $2{,}610$     & $0.644 \pm 0.044$ & $0.688 \pm 0.054$ & $0.628 \pm 0.038$ & $0.465 \pm 0.004$ \\
CNN & joint  & $64$   & $10{,}434$    & $0.574 \pm 0.038$ & $0.703 \pm 0.021$ & $0.698 \pm 0.035$ & $0.421 \pm 0.004$ \\
CNN & frozen & $64$   & $10{,}434$    & $0.611 \pm 0.042$ & $0.688 \pm 0.054$ & $0.658 \pm 0.036$ & $0.437 \pm 0.004$ \\
CNN & joint  & $256$  & $41{,}730$    & $0.540 \pm 0.023$ & $0.712 \pm 0.006$ & $1.070 \pm 0.112$ & $0.290 \pm 0.009$ \\
CNN & frozen & $256$  & $41{,}730$    & $0.574 \pm 0.026$ & $0.690 \pm 0.053$ & $0.787 \pm 0.029$ & $0.345 \pm 0.008$ \\
CNN & joint  & $1024$ & $166{,}914$   & $0.534 \pm 0.013$ & $0.711 \pm 0.007$ & $2.334 \pm 0.169$ & $0.029 \pm 0.004$ \\
CNN & frozen & $1024$ & $166{,}914$   & $0.560 \pm 0.024$ & $0.690 \pm 0.051$ & $1.119 \pm 0.054$ & $0.153 \pm 0.019$ \\
\midrule
ViT & joint  & $64$   & $84{,}546$    & $0.533 \pm 0.021$ & $0.707 \pm 0.008$ & $3.121 \pm 0.131$ & $0.136 \pm 0.030$ \\
ViT & frozen & $64$   & $84{,}546$    & $0.547 \pm 0.050$ & $0.702 \pm 0.037$ & $3.387 \pm 0.233$ & $0.132 \pm 0.024$ \\
ViT & joint  & $128$  & $300{,}162$   & $0.546 \pm 0.047$ & $0.706 \pm 0.019$ & $3.149 \pm 0.936$ & $0.039 \pm 0.016$ \\
ViT & frozen & $128$  & $300{,}162$   & $0.567 \pm 0.038$ & $0.679 \pm 0.057$ & $3.302 \pm 0.451$ & $0.025 \pm 0.010$ \\
ViT & joint  & $256$  & $1{,}124{,}610$ & $0.503 \pm 0.050$ & $0.707 \pm 0.023$ & $3.779 \pm 0.547$ & $0.031 \pm 0.020$ \\
ViT & frozen & $256$  & $1{,}124{,}610$ & $0.539 \pm 0.007$ & $0.700 \pm 0.035$ & $3.522 \pm 0.160$ & $0.015 \pm 0.001$ \\
\bottomrule
\end{tabular}
\end{table}

Three observations, all descriptive at $n = 3$. The frozen arm's endpoint
accuracy is ahead of the joint arm's at every width and both heads, by
$+0.014$ to $+0.038$; every one of those gaps is smaller than the sum of the
two seed standard deviations, so none is resolvable here. Peak accuracies are
indistinguishable between arms (the largest difference is $0.028$, at ViT
$D = 128$, in favour of the joint arm). The endpoint test loss of the joint
CNN arm diverges with width, reaching $2.334 \pm 0.169$ at $D = 1024$ against
$1.119 \pm 0.054$ frozen, while its training loss falls to $0.029$; the ViT
runs are in a memorization regime at every width, with training losses of
$0.015$--$0.14$ and test losses above $3$, so gradient-ratio observables are
uninformative there.

\subsection{Experiment 1.2 (v5): residual overlap with the fast kernel subspace}
\label{app:exp12v5}

\paragraph{Model and protocol.} This experiment combines the model of the
verified shallow instance with the training protocol of
Appendix~\ref{app:exp12v4}. The model is the linear per-pixel map
$\R^{64} \to \R^{16}$ followed by a two-layer ReLU head of width $D$ under
theorem initialization with $\sigma_b = 0.5$; that head initialization is
required by the verified instance and is not PyTorch's default. Optimizer,
data, loss and schedule are exactly those of Appendix~\ref{app:exp12v4}, with
$D \in \{128, 512, 2048\}$, seeds $42$--$44$ and both a \texttt{joint} and a
\texttt{frozen} arm. The frozen arm here keeps \texttt{requires\_grad} true on
the spectral parameters but excludes them from the optimizer's parameter list;
the trajectory is identical to a hard freeze, and the logged spectral gradient
is the counterfactual signal a frozen encoder would have received. Artifacts:
\texttt{results/exp1\_2v5\_residual\_export.csv} ($198$ rows) and
\texttt{results/exp1\_2v5\_residual\_export\_REPORT.md}.

\paragraph{Probe and normalization.} All kernel quantities are computed on a
fixed probe of $N = 256$ supervised sites with $C = 2$ classes, so the probe
logit space is $\R^{512}$. The probe is drawn once per run with a fixed seed
from the first $16$ training images and held fixed for the whole trajectory;
non-probe pixels are set to the ignore index, so the GGN machinery sees the
same site set. With $z = N^{-1/2}\,\mathrm{stack}(\hat y)$ and averaged
cross-entropy, $r = N^{-1/2}(\mathrm{softmax}(z) - \mathrm{onehot}(y))$,
$J_b = N^{-1/2}\,\partial \hat y/\partial b$ and $K_b = J_bJ_b^\top$, so
$\nabla_b\loss = J_b^\top r$ holds exactly. We write
$a_{\mathrm{eff}} = r^\top K_\thetap r / \norm{r}^2$ and
$\kappa_{\mathrm{eff}} = r^\top K_\phip r / \norm{r}^2$ for the realised
Rayleigh quotients, $f_k$ for the fraction of $\norm{r}^2$ inside the top-$k$
eigenspace of $K_\phip$, and $\mathrm{share} =
a_{\mathrm{eff}}/(a_{\mathrm{eff}} + \kappa_{\mathrm{eff}})$, which coincides
by construction with the instantaneous gradient-flow attribution bound.
Table~\ref{tab:app_v5_ident} gives the three identities the runner
recomputes at each of the $198$ measurements.

\begin{table}[htbp]
\centering
\caption{Experiment 1.2 (v5): validation identities, maximum relative error
over all $198$ measurements, from
\texttt{results/exp1\_2v5\_residual\_export\_REPORT.md} \S2. All three are at
float64 round-off, so the kernels here are the same objects the GGN machinery
measures, in the same normalization.}
\label{tab:app_v5_ident}
\begin{tabular}{lr}
\toprule
Identity & Max.\ relative error \\
\midrule
$a_{\mathrm{eff}}\norm{r}^2 = \norm{\nabla_\thetap\loss}^2$ (autograd) & $1.30 \times 10^{-15}$ \\
$\kappa_{\mathrm{eff}}\norm{r}^2 = \norm{\nabla_\phip\loss}^2$ (autograd) & $1.62 \times 10^{-14}$ \\
$\lambda_{\max}(H^{1/2}K_\phip H^{1/2}) = \lambda_{\max}(G_{\phip\phip})$ (power iteration) & $2.47 \times 10^{-15}$ \\
\bottomrule
\end{tabular}
\end{table}

\paragraph{Headline table.} Table~\ref{tab:app_v5_head} reproduces the joint
arm of \S3 of the report at four log-spaced checkpoints for every width and
seed. The worst-case-coercivity proxy
$\lambda_{\max}(K_\thetap)/(\lambda_{\max}(K_\thetap)+\lambda_{\min}(K_\phip))$
is between $0.9998$ and $1.0000$ at every one of the $198$ rows and is omitted
from the table; $\lambda_{\min}(K_\phip)$ ranges from $2.0 \times 10^{-6}$ to
$1.8 \times 10^{-4}$ and is likewise omitted. Full trajectories at all eleven
checkpoints for all eighteen runs, and the frozen arm, are in the report.

\begin{table}[htbp]
\centering
\footnotesize
\setlength{\tabcolsep}{2.5pt}
\caption{Experiment 1.2 (v5), joint arm: probe measurements at four
checkpoints, every width and seed, from
\texttt{results/exp1\_2v5\_residual\_export\_REPORT.md} \S3. \emph{share} is
$a_{\mathrm{eff}}/(a_{\mathrm{eff}}+\kappa_{\mathrm{eff}})$; \emph{proxy} is
$\lambda_{\max}(K_\thetap)/(\lambda_{\max}(K_\thetap)+\lambda_{\max}(K_\phip))$.
These are probe quantities at eleven checkpoints, not a certificate about the
intervening steps.}
\label{tab:app_v5_head}
\begin{tabular}{rrrrrrrrrrrr}
\toprule
$D$ & seed & step & $a_{\mathrm{eff}}$ & $\kappa_{\mathrm{eff}}$ & $\lambda_{\max}(K_\thetap)$ & $\lambda_{\max}(K_\phip)$ & $f_1$ & $f_{20}$ & $f_{100}$ & share & proxy \\
\midrule
$128$ & $42$ & $0$    & $0.704$ & $0.858$   & $2.393$  & $19.839$  & $1.80\cdot10^{-2}$ & $0.3441$ & $0.4356$ & $0.4510$ & $0.1076$ \\
$128$ & $42$ & $32$   & $0.017$ & $0.011$   & $4.887$  & $21.590$  & $1.48\cdot10^{-5}$ & $0.0390$ & $0.1223$ & $0.6022$ & $0.1846$ \\
$128$ & $42$ & $426$  & $0.012$ & $0.022$   & $3.764$  & $21.884$  & $9.52\cdot10^{-6}$ & $0.0594$ & $0.1296$ & $0.3513$ & $0.1468$ \\
$128$ & $42$ & $2400$ & $0.025$ & $0.043$   & $6.965$  & $22.678$  & $3.68\cdot10^{-6}$ & $0.0288$ & $0.1778$ & $0.3688$ & $0.2350$ \\
$128$ & $43$ & $0$    & $0.299$ & $4.219$   & $1.532$  & $21.735$  & $1.34\cdot10^{-1}$ & $0.4465$ & $0.5173$ & $0.0662$ & $0.0659$ \\
$128$ & $43$ & $32$   & $0.013$ & $0.060$   & $3.275$  & $23.963$  & $1.10\cdot10^{-3}$ & $0.0375$ & $0.1516$ & $0.1755$ & $0.1202$ \\
$128$ & $43$ & $426$  & $0.009$ & $0.018$   & $2.229$  & $24.267$  & $1.81\cdot10^{-4}$ & $0.0298$ & $0.1404$ & $0.3405$ & $0.0841$ \\
$128$ & $43$ & $2400$ & $0.036$ & $0.037$   & $8.669$  & $26.355$  & $2.49\cdot10^{-5}$ & $0.0525$ & $0.1913$ & $0.4927$ & $0.2475$ \\
$128$ & $44$ & $0$    & $0.711$ & $0.948$   & $3.793$  & $16.447$  & $5.30\cdot10^{-3}$ & $0.3073$ & $0.3813$ & $0.4285$ & $0.1874$ \\
$128$ & $44$ & $32$   & $0.049$ & $0.078$   & $4.917$  & $17.693$  & $1.53\cdot10^{-4}$ & $0.0400$ & $0.1711$ & $0.3826$ & $0.2175$ \\
$128$ & $44$ & $426$  & $0.022$ & $0.026$   & $5.010$  & $17.847$  & $8.86\cdot10^{-6}$ & $0.0307$ & $0.1702$ & $0.4611$ & $0.2192$ \\
$128$ & $44$ & $2400$ & $0.029$ & $0.042$   & $6.428$  & $18.976$  & $1.04\cdot10^{-4}$ & $0.0309$ & $0.2020$ & $0.4073$ & $0.2530$ \\
\midrule
$512$ & $42$ & $0$    & $1.098$ & $4.107$   & $2.208$  & $70.380$  & $1.28\cdot10^{-2}$ & $0.5162$ & $0.5947$ & $0.2109$ & $0.0304$ \\
$512$ & $42$ & $32$   & $0.027$ & $0.071$   & $8.021$  & $72.767$  & $1.22\cdot10^{-4}$ & $0.0513$ & $0.1544$ & $0.2763$ & $0.0993$ \\
$512$ & $42$ & $426$  & $0.021$ & $0.041$   & $7.199$  & $72.745$  & $2.62\cdot10^{-5}$ & $0.0402$ & $0.1718$ & $0.3393$ & $0.0901$ \\
$512$ & $42$ & $2400$ & $0.065$ & $0.094$   & $24.622$ & $72.428$  & $4.01\cdot10^{-4}$ & $0.0154$ & $0.1741$ & $0.4101$ & $0.2537$ \\
$512$ & $43$ & $0$    & $0.464$ & $7.095$   & $1.717$  & $67.270$  & $3.67\cdot10^{-2}$ & $0.4456$ & $0.5469$ & $0.0614$ & $0.0249$ \\
$512$ & $43$ & $32$   & $0.023$ & $0.052$   & $4.829$  & $69.680$  & $1.14\cdot10^{-6}$ & $0.0334$ & $0.1328$ & $0.3050$ & $0.0648$ \\
$512$ & $43$ & $426$  & $0.013$ & $0.307$   & $3.550$  & $70.279$  & $2.90\cdot10^{-3}$ & $0.0369$ & $0.1388$ & $0.0413$ & $0.0481$ \\
$512$ & $43$ & $2400$ & $0.077$ & $0.085$   & $8.980$  & $70.889$  & $9.68\cdot10^{-6}$ & $0.0374$ & $0.2226$ & $0.4732$ & $0.1124$ \\
$512$ & $44$ & $0$    & $0.218$ & $14.482$  & $1.749$  & $61.928$  & $2.33\cdot10^{-2}$ & $0.4084$ & $0.4749$ & $0.0148$ & $0.0275$ \\
$512$ & $44$ & $32$   & $0.050$ & $0.236$   & $5.380$  & $64.643$  & $2.66\cdot10^{-4}$ & $0.0460$ & $0.1826$ & $0.1734$ & $0.0768$ \\
$512$ & $44$ & $426$  & $0.019$ & $0.064$   & $4.804$  & $65.177$  & $2.64\cdot10^{-6}$ & $0.0253$ & $0.1555$ & $0.2328$ & $0.0686$ \\
$512$ & $44$ & $2400$ & $0.073$ & $0.113$   & $9.265$  & $68.294$  & $1.12\cdot10^{-5}$ & $0.0144$ & $0.1853$ & $0.3935$ & $0.1195$ \\
\midrule
$2048$ & $42$ & $0$    & $0.525$ & $6.973$   & $1.805$  & $253.786$ & $8.38\cdot10^{-6}$ & $0.3850$ & $0.4677$ & $0.0701$ & $0.0071$ \\
$2048$ & $42$ & $32$   & $0.044$ & $0.324$   & $10.807$ & $254.624$ & $9.03\cdot10^{-4}$ & $0.0373$ & $0.1713$ & $0.1197$ & $0.0407$ \\
$2048$ & $42$ & $426$  & $0.026$ & $0.361$   & $11.223$ & $255.856$ & $9.58\cdot10^{-4}$ & $0.0385$ & $0.1612$ & $0.0663$ & $0.0420$ \\
$2048$ & $42$ & $2400$ & $0.196$ & $0.320$   & $41.018$ & $252.457$ & $2.56\cdot10^{-7}$ & $0.0206$ & $0.1848$ & $0.3796$ & $0.1398$ \\
$2048$ & $43$ & $0$    & $0.179$ & $148.674$ & $1.579$  & $276.775$ & $2.17\cdot10^{-1}$ & $0.6836$ & $0.7327$ & $0.0012$ & $0.0057$ \\
$2048$ & $43$ & $32$   & $0.049$ & $0.354$   & $9.486$  & $274.651$ & $4.62\cdot10^{-4}$ & $0.0193$ & $0.1151$ & $0.1217$ & $0.0334$ \\
$2048$ & $43$ & $426$  & $0.022$ & $1.213$   & $5.681$  & $276.344$ & $3.30\cdot10^{-3}$ & $0.0270$ & $0.1585$ & $0.0175$ & $0.0201$ \\
$2048$ & $43$ & $2400$ & $0.153$ & $0.179$   & $17.424$ & $281.103$ & $8.36\cdot10^{-6}$ & $0.0299$ & $0.1326$ & $0.4617$ & $0.0584$ \\
$2048$ & $44$ & $0$    & $0.495$ & $63.759$  & $2.389$  & $253.499$ & $4.96\cdot10^{-2}$ & $0.4682$ & $0.5447$ & $0.0077$ & $0.0093$ \\
$2048$ & $44$ & $32$   & $0.018$ & $0.257$   & $8.093$  & $252.495$ & $1.78\cdot10^{-4}$ & $0.0237$ & $0.1879$ & $0.0648$ & $0.0311$ \\
$2048$ & $44$ & $426$  & $0.034$ & $1.025$   & $8.764$  & $257.769$ & $1.40\cdot10^{-3}$ & $0.0264$ & $0.1860$ & $0.0323$ & $0.0329$ \\
$2048$ & $44$ & $2400$ & $0.131$ & $0.576$   & $30.982$ & $257.813$ & $1.10\cdot10^{-3}$ & $0.0263$ & $0.1584$ & $0.1848$ & $0.1073$ \\
\bottomrule
\end{tabular}
\end{table}

\paragraph{Cumulative gradient-energy fractions.} Table~\ref{tab:app_v5_cum}
gives $\sum_t \norm{\nabla_\thetap\loss}^2 / \sum_t (\norm{\nabla_\thetap
\loss}^2 + \norm{\nabla_\phip\loss}^2)$ over the whole $2400$-step run, from
the per-step logs the runner writes.

\begin{table}[htbp]
\centering
\caption{Experiment 1.2 (v5): cumulative gradient-energy fraction over the
full run, from \texttt{results/exp1\_2v5\_residual\_export\_REPORT.md} \S6.
These are mini-batch gradient norms over the $32$-image training batches
($8192$ sites), not probe quantities, and the optimizer is Adam, so this is a
gradient-energy attribution and not a decomposition of the realised loss
decrease. On the frozen arm the spectral parameters never move, so that column
is a counterfactual ratio.}
\label{tab:app_v5_cum}
\begin{tabular}{rrrr}
\toprule
$D$ & seed & joint & frozen \\
\midrule
$128$  & $42$ & $0.2896$ & $0.7256$ \\
$128$  & $43$ & $0.0806$ & $0.1883$ \\
$128$  & $44$ & $0.3432$ & $0.6091$ \\
$512$  & $42$ & $0.1303$ & $0.3308$ \\
$512$  & $43$ & $0.0618$ & $0.1230$ \\
$512$  & $44$ & $0.0415$ & $0.1336$ \\
$2048$ & $42$ & $0.0173$ & $0.0642$ \\
$2048$ & $43$ & $0.0042$ & $0.0059$ \\
$2048$ & $44$ & $0.0124$ & $0.0310$ \\
\bottomrule
\end{tabular}
\end{table}

\paragraph{What the numbers say.} The following are the report's own
summaries, over the joint arm's $99$ rows unless stated otherwise. The spatial
kernel is extremely ill-conditioned on the probe logit space:
$\lambda_{\max}/\lambda_{\min}$ has median $1.28 \times 10^{6}$ (range
$2.42\times10^{5}$ to $1.06\times10^{7}$), so a coercivity hypothesis
$K_\phip \succeq \kappa I$ taken on the whole output space is available only
with $\kappa = \lambda_{\min}$ and yields nothing. The visited residual is not
in that near-null part --- $\kappa_{\mathrm{eff}}$ exceeds $\lambda_{\min}$ by
between $225\times$ and a median $5.0 \times 10^{3}\times$ --- and it is not in
the coercive top either: $\kappa_{\mathrm{eff}}$ is a median $0.0025$ of
$\lambda_{\max}$ (range $2.5\times10^{-4}$ to $0.537$), and after step $75$ the
top eigendirection carries a median $8.6\times10^{-5}$ of $\norm{r}^2$ (down
from a median $0.0233$ at $t=0$), the top-$20$ eigenspace a median $0.0298$
(down from $0.446$) and the top-$100$ of $512$ a median $0.1612$. The spectral
side collapses the same way: the top left singular direction of $J_\thetap$
carries a median $6.2\times10^{-4}$ of $\norm{r}^2$ after step $75$, down from
$0.1513$ at $t=0$. Measured against the realised share, the top-eigenvalue
proxy understates the encoder's share in $85\%$ of joint rows, by a median
factor $2.72$ and up to $11.6\times$, and overstates it by up to $11.1\times$
in the remaining $15\%$; it errs in both directions and bounds the share in
neither. This is how we read those rows: the diagnostic limitation illustrated
by the counterexample to top-eigenvalue certificates is observed on this
instance.

\paragraph{Scope of these measurements.} Four restrictions apply and are not
optional caveats. (i) The cumulative column is a \emph{gradient-energy
fraction} on the same sequence, not an attribution of Adam's loss decrease;
the algebraic weighted-sum identity holds for any optimizer on one sequence,
but the realised displacement under Adam is not proportional to the gradient.
(ii) Quantities (i)--(v) of the report are computed on the fixed $256$-site
probe while the cumulative column uses training mini-batches: these are the
same functional on different sequences, and we report probe results as probe
results. Eleven checkpoints cannot certify an almost-everywhere trajectory
condition or the intervening steps. (iii) The appropriate isotropic reference
for the subspace fractions lives in the per-site class-contrast space: with
$\Pi = I_N \otimes (I_C - \mathbf{1}\mathbf{1}^\top/C)$, a random unit residual
there has expected projected mass $\mathrm{tr}(P\Pi)/[N(C-1)]$ rather than
$\mathrm{rank}(P)/(NC)$, and for $C = 2$ the difference is consequential. That
reference has not been recomputed for the projectors actually used, so we make
no claim that any measured mass falls below an isotropic level; the measured
masses themselves stand. (iv) The cross-entropy Hessian weighting is a
different quantity, not a correction to a squared-loss approximation: for
cross-entropy $r^\top K_b r = \norm{\nabla_b\loss}^2$ already holds exactly.
For completeness, \S5 of the report records that the weighted share moves by
at most a factor $2.721$ and at least $0.348$ relative to the unweighted one on
this instance. Finally, a cumulative $O(1/D)$ law does not follow from these
three widths: under the phase-restricted attribution result of the mathematics
appendix, a total share of order $1/D$ requires the complementary phase to
carry an $O(1/D)$ fraction of the energy, and that fraction has not been
measured here.

\subsection{Experiment 1.3: the equal-information calibration}
\label{app:exp13}

\paragraph{Calibration.} ``Equal information'' between the spectral pathway
(per-pixel content of $X$ along a direction $u$) and the spatial pathway
(content along an orthogonal direction $v \perp u$, modulated by a smooth
per-position function of the class) requires an operational definition. Three
were compared: \texttt{bayes}, which sets $\alpha, \beta$ so that a
single-modality logistic regression on $\alpha$-only data and a
position-majority classifier on $\beta$-only data both reach a target accuracy
of $0.75$; \texttt{ntk}, which equalizes the leading Gram-matrix eigenvalues of
the two modalities' per-pixel features; and \texttt{margin}, which equalizes
the $L_2$ margins of the two single-modality logistic regressions.

\paragraph{Setup and results.} For each mode the CNN composition is trained at
$D = 256$ for $150$ epochs on $N_{\mathrm{train}} = 512$ samples over three
seeds, in four conditions: joint, frozen at random spectral initialization,
spectral-only ($\beta = 0$) and spatial-only ($\alpha = 0$).
Table~\ref{tab:app_calib} compares the three modes.

\begin{table}[htbp]
\centering
\caption{Experiment 1.3: calibration-mode comparison (CNN, $D = 256$, three
seeds). \emph{spec\,/\,spat} are the achievable single-modality accuracies;
\emph{gap} is $a_{\mathrm{frozen}} - a_{\mathrm{joint}}$.}
\label{tab:app_calib}
\begin{tabular}{lccccc}
\toprule
Mode & $\alpha$, $\beta$ & spec\,/\,spat & joint & frozen & gap \\
\midrule
\texttt{bayes}  & $0.82,\ 18.75$ & $0.75\,/\,0.75$ & $0.48$ & $0.61$ & $+0.13$ \\
\texttt{ntk}    & $0.99,\ 4.91$  & $0.78\,/\,0.58$ & $0.53$ & $0.57$ & $+0.04$ \\
\texttt{margin} & $5.00,\ 5.00$  & $0.95\,/\,0.59$ & $0.60$ & $0.69$ & $+0.09$ \\
\bottomrule
\end{tabular}
\end{table}

Under \texttt{bayes} calibration the single-modality baselines are equal by
construction ($\alpha = 0.82$, $\beta = 18.75$ give spectral-only $0.746$ and
spatial-only $0.754$ under the logistic-regression and position-majority
bounds). On the full CNN model, spectral-only training reaches $0.53$ and
spatial-only training reaches $0.53$; frozen-spectral training on the full data
reaches $0.61$; joint training reaches $0.48$.

\paragraph{Two caveats bound this result.} First, the effect size is modest in
absolute terms. The quoted per-condition standard error of roughly $0.02$ is a
pooled figure over three seeds and $128$ test samples and does not include
between-seed variance, so it understates the uncertainty; at three seeds the
intervals are loose, and $0.48$ is not distinguishable from the $0.50$ chance
floor. Second, the CNN reaches only $0.53$ on single-modality data against a
$0.75$ calibration target, so the network is operating far from the Bayes
bound and part of the joint-training deficit could in principle reflect
generic overfitting with more trainable parameters. What the direction of the
result establishes is that joint training here undershoots baselines that
demonstrably use only one modality, which is not what a reading in which the
spectral features are simply irrelevant would produce.

\subsection{Experiment 1.6: comparison with a logit-magnitude penalty}
\label{app:exp16}

\paragraph{Design.} $120$ runs at width $256$ for $100$ epochs with weight
decay $0$, sweeping the conditions \texttt{joint}, \texttt{frozen} and
\texttt{sd} (spectral decoupling, $(\lambda/2)\norm{\text{logits}}^2$ on the
pre-softmax logits, summed over classes and averaged over examples, following
\citet{pezeshki2021gradient}) with $\lambda \in \{10^{-3}, 10^{-2}, 10^{-1},
1\}$, over noise levels $\{0.05, 0.10, 0.15, 0.20\}$ and five seeds.
Artifacts: \texttt{results/exp1\_6\_summary.csv} ($120$ rows),
\texttt{results/exp1\_6\_notes.md} and per-run logs under
\texttt{results/exp1\_6/}.

\paragraph{Results.} Table~\ref{tab:app_exp16} aggregates over the five seeds
and four noise levels; the means were computed by us from the rows of
\texttt{results/exp1\_6\_summary.csv}, with $\lambda$ selected per noise level
by highest mean final test accuracy. Paired Wilcoxon tests matched on
$(\text{noise}, \text{seed})$, $n = 20$, recomputed from the same rows, give
$T = 88$, $p = 0.55$ for joint against the selected-$\lambda$ penalty, and
$T = 6$, $p = 2.7\times10^{-5}$ for joint against frozen.

\begin{table}[htbp]
\centering
\caption{Experiment 1.6: peak and final test accuracy and their difference,
averaged over five seeds and four noise levels, computed from
\texttt{results/exp1\_6\_summary.csv}. The penalty row uses the $\lambda$ with
the highest mean final accuracy at each noise level ($10^{-2}$ at noise
$0.05$, $10^{-1}$ at the other three).}
\label{tab:app_exp16}
\begin{tabular}{lccc}
\toprule
Condition & Peak acc. & Final acc. & Peak $-$ final \\
\midrule
joint                       & $0.658$ & $0.548$ & $0.111$ \\
logit penalty, selected $\lambda$ & $0.660$ & $0.550$ & $0.110$ \\
frozen                      & $0.583$ & $0.517$ & $0.066$ \\
\bottomrule
\end{tabular}
\end{table}

\paragraph{Caveats.} This experiment was asserted in an earlier draft of the
work without being presented; it is set out here in full, and the following
qualifications belong with it. (i) The frozen condition trades peak accuracy
for a smaller peak-to-final difference: its mean peak is $0.583$ against
$0.660$, because a random spectral projection sometimes destroys signal.
(ii) Its seed-to-seed variability is large --- the standard deviation of final
accuracy across seeds is $0.085$--$0.141$ for frozen against $0.018$--$0.040$
for the other two conditions --- and one seed at noise $0.20$ ends at final
accuracy $0.372$, below chance. (iii) At width $256$ the joint condition does
not collapse dramatically ($0.66$ peak to $0.55$ final); larger drops appear at
$D = 2048$ in Experiment 1.2, so this is a moderate-capacity comparison.
(iv) The best penalty strength depends on the noise level, so a single
$\lambda$ across noise levels is not the best available setting for that
intervention. We report the comparison descriptively, as a contrast between
two interventions on one synthetic instance; nothing here is evaluated as a
recommendation for practice.

\subsection{Solvable serial instance: numerical verification}
\label{app:toy}

\paragraph{Zero-context instance.} Table~\ref{tab:app_toy} verifies the closed
forms and the bounds of the solvable serial cross-entropy instance at
$a_0 = 0$, $v_0 = 1$ and margin $m = \log 19 = 2.944$, from
\texttt{results/toy\_serial\_ce.csv}. The solver column is a stiff ODE
integration of the full $(2+M)$-dimensional state; a separate explicit-Euler
gradient-descent path, run at four step sizes, has observed convergence order
$1.0000004$ and agrees with an independent numpy analytic-gradient
implementation to $0$ relative error.

\begin{table}[htbp]
\centering
\small
\caption{Solvable serial instance ($a_0 = 0$, $v_0 = 1$, $m = \log 19$):
closed form against ODE solver, from \texttt{results/toy\_serial\_ce.csv}. The
last block is the two-timescale displacement bound against the realised
encoder displacement. Every relative discrepancy between closed form and
solver in this table is at most $1.2\times10^{-13}$.}
\label{tab:app_toy}
\begin{tabular}{lrrrr}
\toprule
Quantity & $M = 16$ & $M = 64$ & $M = 256$ & $M = 1024$ \\
\midrule
$\lambda_\thetap(0)$ & $0.25$ & $0.25$ & $0.25$ & $0.25$ \\
$\lambda_\phip(0)$ & $4.0$ & $16.0$ & $64.0$ & $256.0$ \\
$\Dcurv(0)$ & $16.0$ & $64.0$ & $256.0$ & $1024.0$ \\
\midrule
$a(T_m)$ & $0.142325$ & $0.042094$ & $0.011216$ & $0.002857$ \\
suppression $m/a_M$ & $20.688$ & $69.949$ & $262.53$ & $1030.7$ \\
$M\,a_M$ (limit $m$) & $2.2772$ & $2.6940$ & $2.8712$ & $2.9253$ \\
$T_m$ & $0.90784$ & $0.28680$ & $0.078763$ & $0.020252$ \\
$T_m$ lower bound $\Psi(m)/(M{+}1{+}2m^2)$ & $0.60992$ & $0.25437$ & $0.076345$ & $0.020094$ \\
$T_m$ upper bound $\Psi(m)/(M{+}1)$ & $1.2320$ & $0.32222$ & $0.081496$ & $0.020434$ \\
\midrule
$\sup_t\,(q_J - q_F)$ & $0.38232$ & $0.13813$ & $0.039802$ & $0.010376$ \\
uniform gap bound & $11.297$ & $2.9547$ & $0.74729$ & $0.18737$ \\
$(M{+}1)\times$ gap & $6.4994$ & $8.9787$ & $10.229$ & $10.635$ \\
\midrule
$\norm{W(T_m) - W(0)}_2$ & $0.219021$ & $0.071050$ & $0.019658$ & $0.005064$ \\
displacement bound & $1.370601$ & $0.358465$ & $0.090662$ & $0.022732$ \\
bound\,/\,displacement & $6.258$ & $5.045$ & $4.612$ & $4.489$ \\
\midrule
normalized-readout control: $a(T_m)$ & $1.026947$ & $1.026947$ & $1.026947$ & $1.026947$ \\
normalized-readout control: $\norm{W(T_m){-}W(0)}$ & $1.177117$ & $1.177117$ & $1.177117$ & $1.177117$ \\
\midrule
residual envelope excess on $[0, 20T_m]$ & $0$ & $0$ & $0$ & $0$ \\
readout spread $\max_j\beta_j - \min_j\beta_j$ at $T_m$ & $0$ & $0$ & $1.4\cdot10^{-17}$ & $0$ \\
invariant $q = Q_M(a)$, max rel.\ error & $2.0\cdot10^{-11}$ & $6.7\cdot10^{-13}$ & $1.5\cdot10^{-13}$ & $6.8\cdot10^{-14}$ \\
\bottomrule
\end{tabular}
\end{table}

The bound stays within $4.5$ to $6.3$ times the realised displacement across
this width range while both fall by more than a factor of $40$. The
normalized-readout control, in which the readout is $M^{-1/2}Uh$ with $O(1)$
entries trained at unit rate, removes the width dependence exactly: the
encoder's movement is the $M = 1$ value at every width, to $1.7\times10^{-14}$
relative error. The eigenvalue constants $\lambda_\thetap(0) = 1/4$ and
$\lambda_\phip(0) = M/4$ hold for the mean-over-examples binary cross-entropy
with a unit-scaled scalar logit; summing instead of averaging, or rescaling the
logit, multiplies both blocks equally and leaves $\Dcurv(0) = M$ unchanged.

\paragraph{Isotropic extension.} Table~\ref{tab:app_toy_iso} collects the
verification of the general-initialization version, in which $a(0) = a_0$,
$v(0) = v_0 \neq 0$ and the readouts start unequal but sum to exactly zero.
The grid is $M \in \{1, 4, 16, 64, 256, 1024, 4096\}$, $a_0 \in \{-1.5, -0.3,
0, 0.4, 1.7\}$, $v_0 \in \{-1.3, -0.2, 0.05, 0.2, 1.0\}$ and $m \in \{2.0,
2.9, 4.0\}$: $525$ non-trivial cases plus $54$ immediate-stop cases with
$a_0 \geq m$, integrated with DOP853 at \texttt{rtol}~$10^{-10}$ and
\texttt{atol}~$10^{-12}$ with a terminal event at $q = m$.

\begin{table}[htbp]
\centering
\small
\setlength{\tabcolsep}{4pt}
\caption{Isotropic serial instance: verification summary and worst-case bound
slacks over the whole grid, from
\texttt{results/toy\_serial\_ce\_isotropic\_REPORT.md} \S1--\S2. Slacks are
$\text{bound} - \text{value}$ for upper bounds and $\text{value} -
\text{bound}$ for lower bounds; a negative entry would be a violation, and
there are none. $\Delta = m - a_0$.}
\label{tab:app_toy_iso}
\begin{tabular}{lr}
\toprule
Quantity & Value \\
\midrule
Max closed-form vs ODE relative error over $u$, $v(T_m)$, $b(T_m)$, $T_m$ (525 cases) & $1.260\times10^{-11}$ \\
Max relative error of the conserved quantity along the whole trajectory & $1.518\times10^{-10}$ \\
Max root accuracy $|Q_{\mathrm{iso}}(u_M) - m|$ & $3.553\times10^{-15}$ \\
Max $|\sum_j \beta_j(0)|$ (zero-sum initialization is bitwise exact) & $0$ \\
Max spread of $\beta_j(T_m) - \beta_j(0)$ & $7.772\times10^{-16}$ \\
Max frozen-comparator ODE vs closed-form relative error & $5.993\times10^{-11}$ \\
Max $M^{-1/2}$-readout control vs its $M = 1$ value & $1.335\times10^{-12}$ \\
Bound violations, all cases and all bounds & $0$ \\
\midrule
Min slack, $u_M \le \Delta/(1 + Mv_0^2)$ & $2.221\times10^{-10}$ \\
Min slack, $u_M > 0$ & $4.333\times10^{-5}$ \\
Min slack, $|v(T_m) - v_0| \le \Delta^2/(2M|v_0|^3)$ & $1.493\times10^{-9}$ \\
Min slack, encoder displacement bound & $3.918\times10^{-10}$ \\
Min slack, $T_m \le \Psi_{a_0}(m)/(1 + Mv_0^2)$ & $1.748\times10^{-9}$ \\
Min slack, $T_m \ge \Psi_{a_0}(m)/(1 + Mv_0^2 + 2\Delta^2/v_0^2)$ & $3.175\times10^{-9}$ \\
Min slack, $q_J(t) \ge q_F(t)$ & $0$ \\
Min slack, $\sup_t (q_J - q_F) \le C_0/(p_0 M v_0^2)$ & $1.495\times10^{-4}$ \\
\midrule
Reversal error at $\sigma = 1$, $m = 2$, $M = 64\,/\,1024\,/\,16384$ & $0.8207\,/\,0.8380\,/\,0.8397$ \\
Monte-Carlo limit $\Phi(m/2\sigma)$ on the same $4000$ draws & $0.8413$ \\
Draws disagreeing with the limit that satisfy $Mv_0^2 > m/(m - 2a_0)$ & $0$ of $48000$ \\
Spectral-only comparator reversal error, every case & $0$ \\
\bottomrule
\end{tabular}
\end{table}

The update fraction $(a(T_m) - a_0)/(m - a_0)$ falls towards zero in $M$ at
every $(a_0, v_0)$ on the grid --- for instance $0.288$, $0.0423$, $0.000964$
at $M = 1, 16, 1024$ with $a_0 = -1.5$, $v_0 = 1$, $m = 2.9$ --- while
$a(T_m)/m$ and $a(T_m)/a_0$ converge to $a_0/m$ and to $1$ instead, so those
two ratios are not the quantities that vanish. Under the $M^{-1/2}$ readout
the suppression is $0.336399$ at every width, against $0.336$, $0.0472$ and
$0.000969$ for the unnormalized readout at $M = 1$, $16$ and $1024$. Every
Monte-Carlo draw that falls short of the reversal limit fails the sufficient
finite-width condition, either because $|v_0|$ is small or because $a_0$ sits
just below $m/2$, where the threshold diverges.


\section{Production model details}
\label{app:production}

\subsection{Pipeline, capacities and data}
\label{app:prod_setup}

The data are $336 \times 336$ pixel tissue-microarray cores collected by
quantum-cascade-laser infrared spectroscopy over the $1000$--$1800$~cm$^{-1}$
fingerprint region, with per-pixel four-class tissue labels. Each pixel carries
an $S = 942$ dimensional vector, obtained by stacking baseline-corrected
absorbance with its first two derivatives. The compositional model is the Block
Vision Transformer of \citet{oleary2026spatial}: a linear per-pixel
$f_\thetap : \R^{942} \to \R^{K}$, then $16 \times 16$ patch tokenization, a
12-layer, 12-head vision transformer of embedding width $M$, a
transposed-convolution decoder and a convolutional segmentation head. Parameter
counts in Table~\ref{tab:app_capacity} are read from the instantiated model.

\begin{table}[htbp]
\centering
\caption{Module capacities of the production pipeline, counted from the
instantiated model. $K$ is the spectral bottleneck width. The last row is the
single-channel field-standard input convention, for comparison.}
\label{tab:app_capacity}
\begin{tabular}{lrrr}
\toprule
Configuration & $\Cf$ & $\Cg$ & $\Cg/\Cf$ \\
\midrule
$K = 64$, three-channel input  & $61{,}108$  & $18{,}271{,}540$ & $299$ \\
$K = 128$, three-channel input & $121{,}588$ & $21{,}472{,}564$ & $177$ \\
$K = 64$, single-channel input & $20{,}916$  & $18{,}271{,}540$ & $874$ \\
\bottomrule
\end{tabular}
\end{table}

\paragraph{Channel convention.} The published architecture operates on the
second-derivative spectrum alone, the field-standard preprocessing in infrared
histopathology. The models analysed here instead stack raw absorbance with its
first two derivatives, following a companion tokenization study (anonymous,
under review). The choice affects only $\Cf$, and the three-channel convention
is the conservative one for the present purpose: under the field standard the
capacity asymmetry is roughly three times larger. Within the spatial module the
dominant blocks are the transformer ($5.34 \times 10^6$ parameters) and the
transposed convolution ($9.44 \times 10^6$), both $\Theta(M^2)$ in the
embedding width.

\paragraph{The head hypothesis is structurally unmet.} The segmentation head is
the chain $\mathrm{Conv}_{3\times3}(M{+}K \to 96) \to \mathrm{BN} \to
\mathrm{ReLU} \to \mathrm{Conv}_{3\times3}(96 \to 48) \to \mathrm{BN} \to
\mathrm{ReLU} \to \mathrm{Conv}_{3\times3}(48 \to 4)$, so the final affine
classifier sees $48 \times 3 \times 3 = 432$ input features at every width we
build ($M \in \{48, 192, 384\}$, $K \in \{64, 128\}$): the $\Theta(M)$ feature
map is separated from the logits by two nonlinearities. The
curvature-disparity theorem requires a dense classifier over penultimate
features of dimension $\Theta(M)$ with no downstream nonlinearity, and that
hypothesis is not satisfied here. Failing a hypothesis removes a guarantee; it
does not reverse the conclusion.

\subsection{Geometry at initialization}
\label{app:prod_geometry}

Gauss--Newton blocks of the instantiated pipeline were measured on real
training batches at random initialization, by Lanczos with full
reorthogonalization on matrix-free GGN block operators, validated to $10^{-12}$
against dense computation, sweeping $M \in \{48, 96, 192, 384\}$ with three
seeds and four batches per width. Table~\ref{tab:app_geom} collects the
results.

\begin{table}[htbp]
\centering
\small
\setlength{\tabcolsep}{4pt}
\caption{Production geometry at initialization. Ratios are mean $\pm$ sd over
three seeds and four batches per width. Slopes are log--log fits over the width
sweep. Effective ranks are out of $S = 942$. Directional entries are measured
at $M = 192$ over the nine measurements whose batch contains both classes
(three seeds $\times$ three batches).}
\label{tab:app_geom}
\begin{tabular}{lcc}
\toprule
Quantity & Breast & Prostate \\
\midrule
$\Dcurv$ at $M = 48$ & $0.33 \pm 0.12$ & --- \\
$\Dcurv$ at $M = 192$ & $0.46 \pm 0.21$ & --- \\
$\Dcurv$ at $M = 384$ & $0.70 \pm 0.19$ & $1.09 \pm 0.52$ \\
log--log slope of $\lambda_{\max}(G_{\thetap\thetap})$ vs $M$ & $-0.003$ & $-0.012$ \\
log--log slope of $\lambda_{\max}(G_{\phip\phip})$ vs $M$ & $0.38$ ($115 \to 258$) & $0.45$ \\
Effective rank of $\Sigma_X$, raw inputs & $1.07$ & $1.02$ \\
Effective rank of $\Sigma_X$, post-normalization & $1.21$ & --- \\
$|\cos|$ between $v_1$ and the mean spectrum (post-norm.\,/\,raw) & $>0.995$\,/\,$>0.9999$ & --- \\
$\thetap$-block top eigenvector's norm fraction in the $v_1$-induced subspace & $0.954 \pm 0.017$ & --- \\
Trace fraction in the top $40$ of $61{,}108$ spectral directions & $0.83 \pm 0.12$ & --- \\
Curvature ratio, $v_1$ vs full class contrast ($|\cos(d,v_1)| = 0.45$) & $3.7\times$ ($2.6$--$4.8$) & --- \\
Curvature ratio, $v_1$ vs $v_1$-orthogonalized contrast & $12$--$28\times$ & --- \\
Squared-norm fraction of the contrast retained after orthogonalization & $76$--$84\%$ & --- \\
Rank at which $\lambda_{\max}(G_{\phip\phip})$ overtakes the $\thetap$ spectrum & $2$--$5$ & --- \\
\bottomrule
\end{tabular}
\end{table}

Two controls bracket the denominator. Removing the spectral normalization
deflates $\lambda_{\max}(G_{\thetap\thetap})$ by $2.4\times$ ($471 \to 194$ at
$M = 192$) and pushes the ratio across parity ($1.18$ at $M = 192$, $1.31$ at
$M = 384$). Input \emph{centering} changes nothing when the architecture
contains internal batch normalization: $\Dcurv$ is identical to three decimals
with and without centering, because a linear projection followed by
normalization absorbs constant input shifts exactly. The near-rank-one input
statistics make the cap side of the curvature-disparity bound uninformative at
buildable widths independently of the head-hypothesis failure: the data
concentrate almost all their second-moment mass in one direction, and that
direction is the mean spectrum. The directional measurements establish an
anisotropy, not its consequence for learning; whether it starves learning is a
hypothesis, not a measurement.

\subsection{What carries the width trend}
\label{app:prod_width}

The measurements in this subsection are made on the production model at random
initialization, at three widths $M \in \{48, 192, 384\}$ and three seeds, on
one core (the densest of the first $24$ breast fold-0 training cores, with
$29{,}664$ labelled pixels of $112{,}896$), with cuDNN TF32 disabled. Four
normalization modes are compared: \texttt{train} (both head batch
normalizations using batch statistics), \texttt{eval\_bn1} and
\texttt{eval\_bn2} (only the first or only the second head normalization at its
initial running statistics) and \texttt{eval\_head} (both). Artifacts:
\texttt{results/exp1\_8c\_interior\_witness.csv} and
\texttt{results/exp1\_8d\_REPORT.md}.

\paragraph{The incoming features are flat in width.} The only downstream layer
whose input dimension grows with $M$ is the first head convolution. The top
eigenvalue of the patch-unfolded second-moment matrix of the features entering
it is unchanged to four significant figures across the three widths, at
$356.3$, $403.7$ and $332.4$ for seeds $0$, $1$ and $2$ respectively; the
leading direction sits inside the width-independent $K = 64$ skip channels.
Table~\ref{tab:app_gram} also records a discrepancy that matters for any
normalization argument: the same Gram computed over the full batch-normalization
group (all positions, ignored pixels included) rather than over labelled pixels
is substantially smaller, and the cosine between the two centered top
eigenvectors is far from one.

\begin{table}[htbp]
\centering
\small
\caption{Incoming $3\times3$-patch Gram of the first head convolution, top
eigenvalue, from \texttt{results/exp1\_8d\_REPORT.md}. Values are identical
across $M \in \{48, 192, 384\}$ and across all four normalization modes, so a
single column suffices. \emph{labelled} uses the labelled-pixel mask;
\emph{full group} uses the batch-normalization group. The last column is the
cosine between the two centered top eigenvectors.}
\label{tab:app_gram}
\begin{tabular}{rrrrr}
\toprule
seed & labelled & labelled, centered & full group & $\cos$(labelled, full) \\
\midrule
$0$ & $356.3$ & $68.19$ & $263.3$ & $0.2793$ \\
$1$ & $403.7$ & $54.60$ & $306.8$ & $0.5502$ \\
$2$ & $332.4$ & $67.14$ & $253.0$ & $0.2176$ \\
\bottomrule
\end{tabular}
\end{table}

\paragraph{Normalization quantities.} Table~\ref{tab:app_bn} gives the
per-channel batch statistics of the two head normalizations, measured over the
group they actually normalize, together with the activation energy entering and
leaving the first one. Under fan-in initialization the pre-normalization
variance of the first head normalization falls roughly as $1/(M+K)$, and its
inverse-standard-deviation gain grows correspondingly. The stabilizer
$\varepsilon = 10^{-5}$ is never within three orders of magnitude of the
variance: the fraction of channels with variance below $10\varepsilon$ is $0$
in every row of both normalizations at every width, seed and mode.

\begin{table}[htbp]
\centering
\small
\caption{Head normalization quantities at initialization, \texttt{train} mode,
from \texttt{results/exp1\_8d\_REPORT.md} Table 1. Per-seed values are
per-channel medians (variance, gain) or labelled-pixel energies; the
\emph{mean} column is the mean of the three seed values, computed by us.
$\mathrm{gain}_i = \gamma_i/\sqrt{\mathrm{var}_i + \varepsilon}$. The last row
is the product of the mean variance with $M + K$, $K = 64$.}
\label{tab:app_bn}
\begin{tabular}{lrrrrr}
\toprule
Quantity & seed $0$ & seed $1$ & seed $2$ & mean & $M$ \\
\midrule
First head BN, median variance & $0.1470$ & $0.1467$ & $0.1587$ & $0.1508$ & $48$ \\
                               & $0.0678$ & $0.0589$ & $0.0656$ & $0.0641$ & $192$ \\
                               & $0.0346$ & $0.0310$ & $0.0398$ & $0.0351$ & $384$ \\
First head BN, median gain     & $2.609$ & $2.611$ & $2.511$ & $2.577$ & $48$ \\
                               & $3.842$ & $4.120$ & $3.905$ & $3.956$ & $192$ \\
                               & $5.375$ & $5.682$ & $5.010$ & $5.356$ & $384$ \\
Second head BN, median gain    & $2.892$ & $3.122$ & $3.039$ & $3.018$ & $48$ \\
                               & $3.110$ & $3.177$ & $3.195$ & $3.161$ & $192$ \\
                               & $2.966$ & $3.060$ & $3.121$ & $3.049$ & $384$ \\
Pre-BN energy, labelled pixels & $25.64$ & $26.28$ & $23.91$ & $25.28$ & $48$ \\
                               & $11.07$ & $10.67$ & $11.19$ & $10.98$ & $192$ \\
                               & $5.765$ & $5.543$ & $6.015$ & $5.774$ & $384$ \\
Post-BN energy, labelled pixels & $132.1$ & $130.8$ & $133.2$ & $132.0$ & $48$ \\
                                & $130.4$ & $131.5$ & $132.5$ & $131.5$ & $192$ \\
                                & $130.7$ & $128.5$ & $130.8$ & $130.0$ & $384$ \\
\midrule
Mean variance $\times\,(M{+}K)$ & --- & --- & --- & $16.9\,/\,16.4\,/\,15.7$ & $48\,/\,192\,/\,384$ \\
\bottomrule
\end{tabular}
\end{table}

\paragraph{Witness Jacobian-vector products.} A unit perturbation of the first
head convolution is propagated by forward-mode differentiation through the
actual head in the actual normalization mode, so that the train-mode derivative
includes the dependence on the batch statistics. The final scalar
$e_{\mathrm{final}} = N^{-1}\sum_n d_n^\top(\mathrm{diag}(p_n) - p_np_n^\top)d_n$
is a lower-bound witness for $\lambda_{\max}(G_{\phip\phip})$ and is verified
against the same matvec the eigensolver uses; all $216$ rows pass at relative
discrepancy below $10^{-3}$, with a maximum of $1.269 \times 10^{-5}$ over the
$198$ rows with nonzero curvature. Table~\ref{tab:app_jvp} gives the stage
energies for two directions in two modes.

\begin{table}[htbp]
\centering
\small
\caption{Witness JVP stage energies, mean over three seeds, from
\texttt{results/exp1\_8d\_REPORT.md} Table 2a. \texttt{ggn\_top} is the full
top GGN eigenvector of the convolution block; \texttt{s9cen} is the centered
leading direction of the incoming patch Gram. $e_{\mathrm{pre}}$ is the
perturbation energy before the first head normalization,
$e_{\mathrm{logit}}$ is measured on labelled pixels.}
\label{tab:app_jvp}
\begin{tabular}{llrrrrrr}
\toprule
Mode & Direction & $M$ & $e_{\mathrm{pre}}$ & $e_{\mathrm{bn1}}$ & $e_{\mathrm{bn2}}$ & $e_{\mathrm{logit}}$ & $e_{\mathrm{final}}$ \\
\midrule
\texttt{train} & \texttt{ggn\_top} & $48$  & $260.6$ & $2875$ & $4453$ & $261.7$ & $60.73$ \\
\texttt{train} & \texttt{ggn\_top} & $192$ & $259.4$ & $5905$ & $7553$ & $505.2$ & $118.5$ \\
\texttt{train} & \texttt{ggn\_top} & $384$ & $256.8$ & $9168$ & $13810$ & $1018$ & $227.6$ \\
\texttt{train} & \texttt{s9cen}    & $48$  & $76.53$ & $1538$ & $1224$ & $29.23$ & $5.761$ \\
\texttt{train} & \texttt{s9cen}    & $192$ & $76.53$ & $2305$ & $1914$ & $50.83$ & $9.862$ \\
\texttt{train} & \texttt{s9cen}    & $384$ & $76.53$ & $3997$ & $3906$ & $124.3$ & $21.95$ \\
\midrule
\texttt{eval\_head} & \texttt{ggn\_top} & $48$  & $261.9$ & $261.9$ & $42.42$ & $9.607$ & $2.366$ \\
\texttt{eval\_head} & \texttt{ggn\_top} & $192$ & $262.8$ & $262.8$ & $42.18$ & $10.35$ & $2.487$ \\
\texttt{eval\_head} & \texttt{ggn\_top} & $384$ & $262.5$ & $262.5$ & $46.53$ & $10.45$ & $2.372$ \\
\texttt{eval\_head} & \texttt{s9cen}    & $48$  & $76.53$ & $76.53$ & $9.103$ & $0.4667$ & $0.1049$ \\
\texttt{eval\_head} & \texttt{s9cen}    & $192$ & $76.53$ & $76.53$ & $8.460$ & $0.4613$ & $0.1013$ \\
\texttt{eval\_head} & \texttt{s9cen}    & $384$ & $76.53$ & $76.53$ & $9.762$ & $0.4450$ & $0.1011$ \\
\bottomrule
\end{tabular}
\end{table}

A spatially constant bias perturbation $\Delta b = e_c$ is annihilated exactly
by the immediately following train-mode normalization: in \texttt{train} and
\texttt{eval\_bn2} modes, at every width and seed, all stage energies after the
centering step are exactly zero and the quadratic form agrees with the matvec
to at most $1.14\times10^{-17}$ in absolute value. In \texttt{eval\_head} the
same perturbation gives $e_{\mathrm{final}}$ between $7.8\times10^{-4}$ and
$2.0\times10^{-3}$, and in \texttt{eval\_bn1} between $0.0154$ and $0.0692$.

\paragraph{Slopes and attribution.} Table~\ref{tab:app_slopes} gives the
log--log slopes against $M$ per seed and pooled, together with the
mode-to-mode ratios of levels.

\begin{table}[htbp]
\centering
\small
\caption{Log--log slopes against $M$ over the three widths, per seed and
pooled, and mode-to-mode ratios of the levels (each mode's value divided by its
train-mode value at the same width and seed, then averaged over seeds and over
the three widths), from
\texttt{results/exp1\_8d\_REPORT.md}. $\lambda_{WW}$ is the top curvature of
the first head convolution's weight block; $\lambda_\phip$ is the whole spatial
block. Three-width slopes are descriptive, not asymptotic exponents.}
\label{tab:app_slopes}
\begin{tabular}{lrrrrrrr}
\toprule
& \multicolumn{4}{c}{slope of $\lambda_{WW}$} & \multicolumn{1}{c}{$\lambda_\phip$} & \multicolumn{2}{c}{level ratio vs \texttt{train}} \\
\cmidrule(lr){2-5}\cmidrule(lr){6-6}\cmidrule(lr){7-8}
Mode & seed $0$ & seed $1$ & seed $2$ & pooled & pooled slope & $\lambda_{WW}$ & $\lambda_\phip$ \\
\midrule
\texttt{train}     & $0.509$ & $0.785$ & $0.350$ & $+0.613$ & $+0.398$ & $1$ & $1$ \\
\texttt{eval\_bn1} & $0.429$ & $0.634$ & $0.542$ & $+0.563$ & $+0.281$ & $0.903$ & $1.03$ \\
\texttt{eval\_bn2} & $0.544$ & $0.632$ & $0.508$ & $+0.571$ & $+0.253$ & $0.156$ & $0.227$ \\
\texttt{eval\_head} & $-0.112$ & $0.181$ & $-0.084$ & $+0.006$ & $-0.331$ & $0.026$ & $0.037$ \\
\bottomrule
\end{tabular}
\end{table}

Three readings, and one non-reading. The positive width trend in the spatial
block's curvature at initialization is sensitive to head batch normalization:
holding both head normalizations at their initial running statistics removes it
($+0.613 \to +0.006$ for the convolution block, $+0.398 \to -0.331$ for the
whole spatial block) while leaving the incoming features unchanged. Either
normalization alone preserves the trend ($+0.563$ and $+0.571$), which is
evidence against attributing it uniquely to the first one, even though most of
the drop in the \emph{level} attaches to the second. The mechanism this
supports is a finite-range normalization gain acting on width-diluted
activations, over the range of widths measured, and not the width-linear
feature-Gram mechanism of the shallow theory. The non-reading: the
\texttt{eval\_head} intervention changes the normalization derivative, the
forward logits, the downstream gates and the cross-entropy metric at the same
time, so a difference between modes establishes sensitivity to the
intervention, not a unique cause; and the observation that the variance never
approaches the stabilizer is a heuristic about the regime, not a guarantee.

\paragraph{Parameterization control.} Table~\ref{tab:app_scale} scales the
weights and bias of the first head convolution by $c$ and, in the exact
invariance, the following stabilizer by $c^2$. The logits are unchanged and the
block's top curvature scales by exactly $c^{-2}$: the raw Euclidean curvature
of a layer feeding a normalization is a parameterization quantity, not an
architectural one \citep{vanlaarhoven2017l2}.

\begin{table}[htbp]
\centering
\small
\caption{Function-preserving scale control at seed $0$, from
\texttt{results/exp1\_8d\_REPORT.md} Table 3. \emph{ratio} is
$c^2\lambda_{WW}(c)/\lambda_{WW}(1)$, which is $1$ under exact invariance.
\emph{logit dev.} is $\max|\text{logits}(c) - \text{logits}(1)|$, reported both
with the stabilizer co-scaled by $c^2$ and with it held fixed (the
implementation's actual behaviour); $\lambda_{WW}$, \emph{ratio} and
$\norm{W}_F^2$ are taken from the fixed-stabilizer variant, and the co-scaled
variant differs from it by at most one unit in the fourth significant figure.}
\label{tab:app_scale}
\begin{tabular}{rrrrrrr}
\toprule
$M$ & $c$ & $\lambda_{WW}$ & ratio & logit dev.\ (co-scaled) & logit dev.\ (fixed $\varepsilon$) & $\norm{W}_F^2$ \\
\midrule
$48$  & $1/4$ & $719.4$ & $1.000$ & $0$ & $9.40\cdot10^{-4}$ & $2.003$ \\
$48$  & $1/2$ & $179.8$ & $1.000$ & $0$ & $1.89\cdot10^{-4}$ & $8.014$ \\
$48$  & $1$   & $44.96$ & $1.000$ & $0$ & $0$ & $32.06$ \\
$48$  & $2$   & $11.24$ & $1.000$ & $0$ & $4.63\cdot10^{-5}$ & $128.2$ \\
$48$  & $4$   & $2.810$ & $1.000$ & $0$ & $5.84\cdot10^{-5}$ & $512.9$ \\
$192$ & $1/4$ & $1935$  & $0.9998$ & $0$ & $3.00\cdot10^{-3}$ & $2.002$ \\
$192$ & $4$   & $7.560$ & $1.000$ & $0$ & $1.89\cdot10^{-4}$ & $512.6$ \\
$384$ & $1/4$ & $1927$  & $0.9991$ & $0$ & $4.69\cdot10^{-3}$ & $1.999$ \\
$384$ & $4$   & $7.535$ & $1.000$ & $0$ & $2.95\cdot10^{-4}$ & $511.8$ \\
\bottomrule
\end{tabular}
\end{table}

\paragraph{Eigensolver accuracy.} Every row of the width sweep terminates on
the combined criterion (relative change of the Rayleigh quotient below
tolerance \emph{and} residual below $10^{-3}$), never on the iteration or
wall-clock budget. Over all $36$ (mode, width, seed) rows the relative residual
$\norm{Gv - \lambda v}/\lambda$ is at most $3.81\times10^{-4}$ for the
convolution weight block, at most $3.81\times10^{-4}$ for the weight--bias
block and at most $9.97\times10^{-4}$ for the whole spatial block, and the
Rayleigh quotient of the returned iterate agrees with the reported eigenvalue
to the printed precision in every row.

\subsection{Controlled retraining: original protocol}
\label{app:prod_e3c}

The pipeline was retrained on breast fold $0$ at reduced depth (six transformer
layers) with full instrumentation. Two protocols were run and both are
reported: the \emph{original} protocol ($39$ runs), which is the pre-registered
one and whose joint-versus-frozen contrast an audit of the training code later
showed to be confounded, and the \emph{matched} protocol ($21$ runs), which
removes the confounds, saves the selected checkpoint and re-runs the
pre-registered comparisons. Sixty logged runs in total. Per-step residuals,
block gradient norms, per-epoch parameter displacements and input-Jacobian
norms are recorded in both; the matched protocol additionally logs the
per-step clipping coefficient and the per-block applied update norms.

The original protocol comprises arms $\{$joint-linear, frozen-random,
frozen-PCA, joint-MLP$\}$ $\times$ widths $M \in \{48, 192\}$ $\times$ three
seeds under AdamW for $60$ epochs, plus a $\{$joint-linear, frozen-random$\}$
$\times$ three-seed pair under momentum SGD at both widths and a frozen-PCA SGD
triple at $M = 48$ as a positive control. The spectral parameters sit in a
zero-weight-decay group in every arm, and frozen arms log counterfactual
spectral gradients. Table~\ref{tab:app_e3c_orig} gives the AdamW arms.

\begin{table}[htbp]
\centering
\small
\caption{Controlled study, original protocol, AdamW arms: validation macro-F1
(mean $\pm$ sd over three seeds) under two checkpoint policies, and relative
spectral displacement at end of training. \emph{Final-5} is the pre-registered
metric. \emph{Best-val} is the maximum over the $60$ per-epoch evaluations on
the same $28$ validation cores; no checkpoint was saved at that maximum and
this study has no test set, so that column is exploratory and
selection-optimistic. The two policies order the arms oppositely. These arms
also differ in clip scope and in normalization-affine treatment, so this
contrast is confounded.}
\label{tab:app_e3c_orig}
\begin{tabular}{lccccr}
\toprule
& \multicolumn{2}{c}{Final-5 (pre-registered)} & \multicolumn{2}{c}{Best-val (exploratory)} & \\
\cmidrule(lr){2-3}\cmidrule(lr){4-5}
Arm & $M{=}48$ & $M{=}192$ & $M{=}48$ & $M{=}192$ & $\norm{\Delta\thetap}/\norm{\thetap_0}$ \\
\midrule
Joint (linear) & $0.538 \pm 0.082$ & $0.646 \pm 0.050$ & $0.873 \pm 0.030$ & $0.858 \pm 0.056$ & $0.54$--$0.64$ \\
Frozen random  & $0.648 \pm 0.063$ & $0.731 \pm 0.023$ & $0.840 \pm 0.025$ & $0.817 \pm 0.017$ & $0$ \\
Frozen PCA     & $0.731 \pm 0.003$ & $0.726 \pm 0.012$ & $0.769 \pm 0.021$ & $0.862 \pm 0.068$ & $0$ \\
Joint (MLP)    & $0.566 \pm 0.070$ & $0.575 \pm 0.184$ & $0.900 \pm 0.000$ & $0.897 \pm 0.001$ & $0.85$--$0.96$ \\
\bottomrule
\end{tabular}
\end{table}

On the pre-registered final-window metric the joint-linear minus
frozen-random gaps are $-0.110$ at $M = 48$ ($90\%$ paired-$t$ interval
$[-0.333, +0.114]$) and $-0.085$ at $M = 192$ ($[-0.185, +0.015]$), neither
resolvable at $n = 3$; against frozen-PCA they are $-0.193$
($[-0.334, -0.052]$) and $-0.081$ ($[-0.175, +0.014]$). Taking each run's
maximum over its $60$ evaluations reverses the ordering: joint minus
frozen-random becomes $+0.033$ and $+0.041$, joint minus frozen-PCA $+0.104$
and $-0.004$. The switch is an arithmetic identity --- the change in every gap
equals the difference in peak-to-final degradation between the two arms ---
and averaged over seeds that degradation is $0.335$ ($M = 48$) and $0.212$
($M = 192$) for joint-linear against $0.192$ and $0.086$ for frozen-random.

\paragraph{Three asymmetries.} First, gradient-norm clipping at $1.0$ was
computed over the optimizer's parameters, that is over $\thetap + \phip$ in the
joint arms but over $\phip$ alone in the frozen arms, so at identical gradients
the frozen arm takes the larger $\phip$ step: the median extra clip factor paid
by the joint arm is $1.10$--$1.12$ under SGD at $M = 48$ and $3.0$--$3.6$ under
AdamW at $M = 192$. Second, the normalization following the spectral projection
contributes $128$ affine parameters that are trained in joint-linear, frozen in
frozen-random and structurally absent in frozen-PCA. Third, no checkpoint was
saved during training, so the best-validation column is a post-hoc curve
maximum that no saved model realizes, taken over $60$ evaluations of the same
$28$ validation cores that define the comparison.

\subsection{Momentum-SGD controls and normalization recalibration}
\label{app:prod_sgd}

The two optimizer families constitute an optimizer-family control; the SGD arm
is not an instantiation of the two-timescale theorem's gradient flow. It uses
momentum $0.9$; weight decay $0.01$ on $\phip$ and $0$ on $\thetap$;
mini-batches of four cores with reshuffling; and global gradient-norm clipping
at $1.0$, which binds on $73$--$89\%$ of steps. Logged gradient norms are
recorded before clipping, so the plotted theory quantities are not the
quantities the trajectory applied. Under momentum SGD the paired joint minus
frozen-random gap is $+0.002$ at $M = 48$ and $+0.011$ at $M = 192$ on
best-validation, with $90\%$ paired-$t$ intervals $[-0.008, +0.011]$ and
$[-0.007, +0.030]$; on the final-window metric the same pairs give $+0.028$ and
$-0.057$, with intervals $[-0.173, +0.228]$ and $[-0.124, +0.010]$ --- gaps
that differ in sign between the widths, with intervals an order of magnitude
wider than the effect they bracket. Both arms lose $0.12$--$0.22$ macro-F1
between their peak and their final window. Descriptively, the spectral share of
the summed squared gradient norms drops from about $0.94$ (AdamW joint) to
about $0.26$ (SGD joint); momentum, clipping, weight decay and adaptivity all
differ between the two families, so no single factor is isolated.

\paragraph{Recalibration of normalization statistics.} On disposable copies of
the final checkpoints, with all weights fixed, the normalization running
statistics were re-estimated on training data only ($29$ batches) and the model
re-evaluated. For this intervention and these checkpoints, the recovered
fraction of the joint arms' peak-to-final drop is $0.17$ at $M = 48$ and
$0.07$ at $M = 192$: recalibration recovers a nonzero part of the degradation
but not most of it. What it does undo is the instability of the endpoint
number: it repairs the two collapsed joint checkpoints ($+0.24$ and $+0.15$
macro-F1) and shrinks the seed standard deviation by roughly $3.7\times$, while
leaving the frozen arms essentially unchanged ($|\Delta| \le 0.02$, and
$\le 0.006$ for frozen-PCA, whose spectral stage carries no normalization).

\subsection{Matched protocol}
\label{app:prod_matched}

The matched protocol equalizes the three asymmetries and re-runs the
pre-registered comparisons at $M = 48$: the reduction-stage normalization
affine parameters are placed in $\phip$ and trained wherever that
normalization exists; the clip norm at $1.0$ is computed over $\phip$ only
under the same clipping rule and scope in every arm; and the best-validation
checkpoint is saved. AdamW at learning rate $10^{-4}$, $60$ epochs, three
seeds, flip and $90^\circ$-rotation augmentation on in all arms. The logged
per-step clip coefficient averages $0.33$--$0.57$ across arms --- the
coefficients and gradients differ by arm, which is an outcome of the matched
rule --- and the logged applied update norms confirm that the spectral share of
the applied update tracks the learning-rate multiplier ($0.02$, $0.17$ and
$1.36$ times the $\phip$ update norm for $\times0.1$, $\times1$ and
$\times10$). Table~\ref{tab:app_e3c_matched} gives all seven arms.

\begin{table}[htbp]
\centering
\small
\caption{Controlled study, matched protocol ($M = 48$, AdamW, three seeds,
mean $\pm$ sd). Identical normalization-affine treatment and the same
$\phip$-only clipping rule and scope in every arm; best-validation checkpoints
saved. \emph{Final-5} is the pre-registered metric; \emph{best-val} is the
maximum over $60$ per-epoch evaluations on the $28$ fold-0 validation cores and
is exploratory and selection-optimistic, since this study has no test set. The
drift column is the relative spectral-parameter displacement at end of
training; its standard deviations were computed by us from
\texttt{results/e3c\_analysis\_runs.csv}. Rows $1$--$2$ and rows $3$--$4$ are
the two matched pairs; see the text.}
\label{tab:app_e3c_matched}
\begin{tabular}{lccc}
\toprule
Arm & Final-5 (pre-registered) & Best-val (exploratory) & $\norm{\Delta\thetap}/\norm{\thetap_0}$ \\
\midrule
Frozen random                   & $0.684 \pm 0.037$ & $0.814 \pm 0.040$ & $0$ \\
Joint (linear)                  & $0.694 \pm 0.056$ & $0.837 \pm 0.029$ & $0.566 \pm 0.026$ \\
Frozen pretrained               & $0.690 \pm 0.033$ & $0.861 \pm 0.022$ & $0$ \\
Fine-tuned pretrained           & $0.710 \pm 0.052$ & $0.857 \pm 0.009$ & $0.020 \pm 0.005$ \\
Joint, spectral lr $\times 0.1$ & $0.626 \pm 0.055$ & $0.841 \pm 0.007$ & $0.083 \pm 0.006$ \\
Joint, spectral lr $\times 10$  & $0.561 \pm 0.114$ & $0.886 \pm 0.037$ & $4.910 \pm 0.402$ \\
Joint, cosine schedule          & $0.660 \pm 0.020$ & $0.825 \pm 0.033$ & $0.359 \pm 0.013$ \\
\bottomrule
\end{tabular}
\end{table}

\paragraph{Two matched pairs, not a single-factor design.} The four core arms
do not all differ in one factor. Joint-linear and frozen-random use a linear
reduction stage with normalization; frozen-pretrained and fine-tuned-pretrained
use a pretrained multilayer spectral encoder without reduction-stage
normalization. The clean freezing comparisons are therefore the two matched
pairs, joint-linear against frozen-random and fine-tuned against
frozen-pretrained. Comparing across the pairs --- frozen-pretrained against
frozen-random, for instance --- changes architecture, initialization and
training history, and validation selection, not only informativeness.

\paragraph{Verdicts.} All comparisons in Table~\ref{tab:app_verdicts} are
paired by seed with $90\%$ $t$-intervals ($t = 2.920$, $n = 3$), conditional on
one validation split, and are descriptive seed summaries rather than a
population or multi-factor causal analysis.

\begin{table}[htbp]
\centering
\small
\setlength{\tabcolsep}{4pt}
\caption{Matched protocol: paired differences with $90\%$ paired-$t$ intervals
($n = 3$). The last two rows compare the two protocols and are reported as
descriptive seed summaries; saving a checkpoint changes what is available to
evaluate, not the final-window trajectory.}
\label{tab:app_verdicts}
\begin{tabular}{lcc}
\toprule
Contrast & Final-5 & Best-val \\
\midrule
Fine-tuned $-$ frozen-pretrained & $+0.021\ [-0.022, +0.063]$ & $-0.004\ [-0.048, +0.039]$ \\
Joint-linear $-$ frozen-random & $+0.011\ [-0.142, +0.163]$ & $+0.023\ [-0.066, +0.112]$ \\
Frozen-pretrained $-$ frozen-random & --- & $+0.047\ [-0.042, +0.136]$ \\
Spectral lr $\times10 - \times1$ & $-0.133\ [-0.419, +0.152]$ & $+0.048\ [-0.016, +0.112]$ \\
Spectral lr $\times0.1 - \times1$ & $-0.068\ [-0.173, +0.037]$ & $+0.004\ [-0.039, +0.046]$ \\
\midrule
Joint/frozen gap: matched $-$ original & $+0.120\ [-0.230, +0.470]$ & $-0.010\ [-0.105, +0.085]$ \\
\bottomrule
\end{tabular}
\end{table}

The pre-registered prediction that fine-tuning a pretrained spectral encoder
falls below leaving it frozen is not supported: neither difference is
resolvable at $n = 3$, and the fine-tuned arm moved $\thetap$ by only $2\%$ of
its initial norm. The logged applied update norms verify nonzero applied
encoder updates under AdamW in that arm; this is a statement about the applied
updates, not an attribution of the loss decrease. The pre-registered
harmful-drift ordering $\times0.1 \ge \times1 \ge \times10$ is not observed: on
best-validation the $\times10$ arm is above $\times1$ in $3/3$ seeds while its
final-window endpoint is the lowest and most variable of any arm at $0.561$.
Three multiplier means do not establish a systematic relation between
displacement and peak score, and the intervals are wide. The cosine schedule
improved neither mean endpoint ($0.825$ best-val and $0.660$ final-5 against
$0.837$ and $0.694$ for the constant learning rate); one schedule that did not
help does not rule out an excessive late learning rate.

\paragraph{Disclosures.} Width $48$ only; $n = 3$; a single validation fold of
$28$ fold-0 cores; no test set. The best-validation column is exploratory and
the final-window column is the pre-specified metric. The pretrained encoders
were selected by pixel validation macro-F1 on the \emph{same} fold-0 validation
cores the runner later scores, which is a selection advantage for both
pretrained arms, and were pretrained with weight decay $0.01$ on $\thetap$.
Pixel pretraining used the unaugmented centre-crop protocol while subsequent
training used random crop plus augmentation; this affects comparisons across
the two pairs, not the frozen-versus-fine-tuned pair that starts from the same
checkpoint. Augmentation is on in all arms. Finally, the summary artifact
\texttt{results/e3c\_analysis\_runs.csv} does not expose per-run configuration
hashes or data-order pairing, so before the pairs are described as completely
matched a compact provenance manifest should accompany it, recording
architecture, initial checkpoint hash, parameter-group membership, clipping and
weight-decay settings, normalization mode and state, augmentation and
data-order seed, scheduler and selection rule; identical seed labels alone do
not establish bit-identical mini-batch order across different model
constructions.

\subsection{What is not certified on this pipeline}
\label{app:prod_scope}

The spatial module's input-Jacobian operator norm grows over training, from
$3.2$ to $157$ on average for joint-linear (up to $414$) and from $1.7$ to
$1.2\times10^{4}$ for joint-MLP, which rejects an initialization-scale
stability approximation for the containment constant of the displacement
theorem, though not the existence of a finite-horizon constant over the
training window. The fitted per-step residual-envelope rate for joint-linear
grows from $\hat\mu = 6.0\times10^{-4}$ ($M = 48$) to $1.9\times10^{-3}$
($M = 192$), a factor $3.2$ for a $4\times$ width increase; the ratio is
invariant to the step-clock-to-flow-time conversion, the absolute values are
not, and two widths do not establish a trend.

The earlier proxy instantiation of the displacement bound is withdrawn. Even
after correcting the time units and using the supremum gain, the proxy exceeds
the measured displacement by roughly $17$--$47\times$ under AdamW and roughly
$450$--$1900\times$ under momentum SGD; the raw per-step residual violates the
fitted envelope on $19$--$28\%$ of joint-linear steps, so the envelope
assumption is verified for the smoothed curve rather than for the logged
quantity. Momentum, gradient clipping and adaptive preconditioning all lie
outside the theorem's hypotheses, and under AdamW the measured displacement
exceeds $\eta\sum_k\norm{\nabla_\thetap\loss_k}$ by more than an order of
magnitude, so no step-size-rescaled flow bound describes that trajectory in
either direction. No certified instantiation of the displacement bound exists
for this pipeline at present.

Finally, the scalar gradient-ratio diagnostic is blind here: it \emph{rises}
during joint AdamW training (joint-linear epoch-1 median $0.6$--$0.8$, rising
to $3$--$5$ by epoch $60$), including in runs whose validation score is
falling. A single global ratio cannot see a direction-wise anisotropy, and the
initialization geometry of Appendix~\ref{app:prod_geometry} is what such a
ratio would have to resolve.
